% !TEX encoding = UTF-8
% ============================================================
% CHƯƠNG 1: TỔNG QUAN - TEMPLATE
% ============================================================
% Hướng dẫn: 
% - Thay thế [NỘI DUNG] bằng nội dung thực tế
% - Xóa các phần không cần thiết
% - Giữ lại cấu trúc và ví dụ
% ============================================================

\chapter{Tổng quan}

% ============================================================
\section{Đặt vấn đề}
% ============================================================

[VĂN BẢN GIỚI THIỆU VẤN ĐỀ CHÍNH: Bối cảnh, tại sao vấn đề này quan trọng, 
những nỗ lực hiện tại và giới hạn của chúng. ]

Trích dẫn ví dụ: \textcite{tac_gia_sach_2024} cho rằng [NỘI DUNG VĂN BẢN].

\subsection{Vấn đề con 1}
[MÔ TẢ CHI TIẾT VẤN ĐỀ CON 1]

\subsection{Vấn đề con 2}
[MÔ TẢ CHI TIẾT VẤN ĐỀ CON 2]

% ============================================================
\section{Mục tiêu và phạm vi}
% ============================================================

\subsection{Mục tiêu chính}
\begin{itemize}
    \item [MỤC TIÊU 1]
    \item [MỤC TIÊU 2]
    \item [MỤC TIÊU 3]
\end{itemize}

\subsection{Phạm vi nghiên cứu}
\begin{itemize}
    \item [PHẠM VI 1]
    \item [PHẠM VI 2]
    \item [PHẠM VI 3]
\end{itemize}

% ============================================================
\section{Tổng quan về Lĩnh vực Chính}
% ============================================================

[VĂN BẢN GIỚI THIỆU LĨ̀NH VỰC CHÍNH]

\subsection{Khái niệm 1}

[MÔ TẢ CHI TIẾT VỀ KHÁI NIỆM 1]

% Ví dụ công thức toán học
Ví dụ công thức:
\begin{equation}
    y = mx + b
    \label{eq:linear}
\end{equation}

Như trong Phương trình \eqref{eq:linear}, [GIẢI THÍCH].

\subsection{Khái niệm 2}

[MÔ TẢ CHI TIẾT VỀ KHÁI NIỆM 2]

% ============================================================
% VÍ DỤ 1: CHÈN HÌNH ẢNH ĐƠN
% ============================================================
\begin{figure}[H]
    \centering
    \includegraphics[width=0.6\textwidth]{images/[TEN_FILE].png}
    \caption{Tiêu đề hình này là gì. Nguồn: \cite{tac_gia_sach_2024}}
    \label{fig:example_figure}
\end{figure}

Như thể hiện trong Hình \ref{fig:example_figure}, [GIẢI THÍCH HÌNH ẢNH].

% ============================================================
% VÍ DỤ 2: CHÈN NHIỀU HÌNH (SUBFIGURES)
% ============================================================
\begin{figure}[H]
    \centering
    \begin{subfigure}{0.45\textwidth}
        \includegraphics[width=\textwidth]{images/[TEN_FILE_1].png}
        \caption{Hình con 1}
        \label{fig:subfig1}
    \end{subfigure}
    \quad
    \begin{subfigure}{0.45\textwidth}
        \includegraphics[width=\textwidth]{images/[TEN_FILE_2].png}
        \caption{Hình con 2}
        \label{fig:subfig2}
    \end{subfigure}
    \caption{So sánh giữa hai phương pháp. Nguồn: \cite{conference_2024}}
    \label{fig:subfigures}
\end{figure}

% ============================================================
% VÍ DỤ 3: CHÈN BẢNG BIỂU
% ============================================================
\subsection{Khái niệm 3}

[MÔ TẢ CHI TIẾT VỀ KHÁI NIỆM 3]

\begin{table}[H]
    \centering
    \caption{Tiêu đề bảng}
    \label{tab:example_table}
    \begin{tabular}{|l|c|r|}
        \hline
        \textbf{Cột 1} & \textbf{Cột 2} & \textbf{Cột 3} \\
        \hline
        Dữ liệu 1 & 10 & 100 \\
        Dữ liệu 2 & 20 & 200 \\
        Dữ liệu 3 & 30 & 300 \\
        \hline
    \end{tabular}
\end{table}

Từ Bảng \ref{tab:example_table}, ta thấy rằng [GIẢI THÍCH BẢNG].

% ============================================================
% VÍ DỤ 4: CHÈN BẢNG PHỨC TẠP
% ============================================================
\begin{table}[H]
    \centering
    \caption{So sánh các phương pháp}
    \label{tab:comparison}
    \begin{tabularx}{\textwidth}{|l|X|X|X|}
        \hline
        \textbf{Tiêu chí} & \textbf{Phương pháp A} & \textbf{Phương pháp B} & \textbf{Phương pháp C} \\
        \hline
        Hiệu suất & Cao & Trung bình & Thấp \\
        Chi phí & Cao & Trung bình & Thấp \\
        Độ phức tạp & Cao & Trung bình & Thấp \\
        \hline
    \end{tabularx}
\end{table}

% ============================================================
% VÍ DỤ 5: CHÈN MÃ NGUỒN (CODE)
% ============================================================
\subsection{Khái niệm 4}

[MÔ TẢ CHI TIẾT VỀ KHÁI NIỆM 4]

Ví dụ mã nguồn Python:
\begin{minted}{python}
def hello_world():
    print("Hello, World!")
    return True

if __name__ == "__main__":
    hello_world()
\end{minted}

Ví dụ mã nguồn JavaScript:
\begin{minted}{javascript}
const greet = (name) => {
    console.log(`Hello, ${name}!`);
};

greet("World");
\end{minted}

% ============================================================
\section{Đóng góp của nghiên cứu}
% ============================================================

\begin{enumerate}
    \item [ĐÓNG GÓP 1: Mô tả chi tiết]
    \item [ĐÓNG GÓP 2: Mô tả chi tiết]
    \item [ĐÓNG GÓP 3: Mô tả chi tiết]
\end{enumerate}

% ============================================================
\section{Cấu trúc của đồ án}
% ============================================================

Đồ án được tổ chức như sau:

\begin{itemize}
    \item \textbf{Chương 1 - Tổng quan:} Giới thiệu vấn đề, mục tiêu và phạm vi nghiên cứu
    \item \textbf{Chương 2 - [Tiêu đề chương 2]:} [Mô tả nội dung chương 2]
    \item \textbf{Chương 3 - [Tiêu đề chương 3]:} [Mô tả nội dung chương 3]
    \item \textbf{Chương 4 - Kết luận:} Tổng kết, đóng góp và hướng phát triển trong tương lai
\end{itemize}
